% ══════════════════════════════════════════════════════════════════════════════
% ACORDEÓN DE R — Laboratorio de R para Ciencia Política, ITAM
% recursos/acordeon-r.tex · Hoja de referencia de una sola página
%
% No es un deck: no usa el preámbulo de beamer. Es un documento de una sola
% página, tamaño carta, tres columnas, con la paleta de estilo/paleta-itam.tex.
% Solo lleva lo que de verdad se usó en las once sesiones del curso.
%
% Se compila con pdflatex desde esta misma carpeta:
%   pdflatex acordeon-r.tex
% ══════════════════════════════════════════════════════════════════════════════

\documentclass[9pt]{extarticle}

% ── Página: carta, márgenes angostos para que quepa una hoja de referencia ────
\usepackage[letterpaper, margin=0.42in, top=0.38in, bottom=0.38in]{geometry}

% ── Tipografía: sin serifas, misma familia que las láminas del curso ──────────
\usepackage[T1]{fontenc}
\usepackage[utf8]{inputenc}
\usepackage{lmodern}
\usepackage{tgheros}
\renewcommand{\familydefault}{\sfdefault}
\usepackage{microtype}

\usepackage{multicol}
\usepackage{enumitem}
\usepackage{xcolor}
% ══════════════════════════════════════════════════════════════════════════════
% PALETA DEL LABORATORIO DE R PARA CIENCIA POLÍTICA — ITAM
% Verde institucional #006847 sobre blanco
%
% ESTE ES EL ÚNICO LUGAR NORMATIVO DONDE VIVEN LOS HEXES.
% Si cambia un color, se cambia aquí y se propaga a mano a los otros dos
% archivos que lo heredan:
%   · estilo/tema_lab.R     → theme_lab() y las escalas de ggplot2
%   · estilo/estilo.scss    → el sitio Quarto
% Los tres tienen que decir lo mismo. Si no, las láminas, los gráficos y el
% sitio dejan de combinar y se nota.
% ══════════════════════════════════════════════════════════════════════════════

\definecolor{verdeITAM}{HTML}{006847}     % acento principal — verde institucional
\definecolor{verdeProfundo}{HTML}{00402C} % títulos y texto enfático
\definecolor{verdeClaro}{HTML}{A8D5BE}    % fondos suaves, resaltados
\definecolor{grisTexto}{HTML}{3D3D3D}     % cuerpo de texto
\definecolor{guinda}{HTML}{7B1113}        % acento secundario, uso escaso

% ── Derivados ────────────────────────────────────────────────────────────────
% No son colores nuevos: son diluciones de los anteriores. Se definen aquí para
% que nadie escriba "verdeClaro!15" a mano en una lámina y el curso termine con
% quince tonos distintos repartidos por ahí.

\colorlet{verdeVelo}{verdeClaro!22!white}   % fondo de bloque y de código
\colorlet{verdeTenue}{verdeClaro!55!white}  % resaltado dentro de un bloque
\colorlet{verdeSombra}{white!30!verdeITAM}  % texto atenuado SOBRE fondo verde
\colorlet{grisVelo}{grisTexto!7!white}      % fondos muy tenues, reglas
\colorlet{grisTenue}{grisTexto!45!white}    % texto terciario, pies
   % único lugar normativo de los hexes del curso
\usepackage{titlesec}
\usepackage{array}

\pagestyle{empty}
\setlength{\parindent}{0pt}
\setlength{\columnsep}{0.28in}
\setlength{\columnseprule}{0.4pt}
\renewcommand{\columnseprulecolor}{\color{verdeClaro}}

% ── Comando para el encabezado de cada tarea ───────────────────────────────────
% Una barra verde compacta con el nombre de la tarea en blanco. Nada de
% sombras ni degradados, igual que en las láminas del curso.
\newcommand{\tarea}[1]{%
  \par\vspace{8pt}%
  \noindent\colorbox{verdeITAM}{\parbox[c][14pt]{\dimexpr\linewidth-2\fboxsep\relax}{\hspace{2pt}\color{white}\bfseries\fontsize{9.3}{10}\selectfont #1}}%
  \par\vspace{4pt}%
}

% ── Formato compacto de lista: código + glosa en una sola línea ───────────────
\setlist[itemize]{
  leftmargin=7pt, itemsep=3.2pt, parsep=0pt, topsep=2pt, partopsep=0pt,
  labelsep=3pt
}
\renewcommand{\labelitemi}{\tiny\textcolor{verdeITAM}{$\bullet$}}

\newcommand{\co}[1]{\texttt{\color{verdeProfundo}#1}}

\begin{document}
\fontsize{9}{10.3}\selectfont

% ── Cabecera de la hoja ─────────────────────────────────────────────────────
\begin{center}
  {\color{verdeProfundo}\bfseries\fontsize{13}{13}\selectfont Acordeón de R}\quad
  {\color{grisTexto}\fontsize{8}{8}\selectfont Laboratorio de R para Ciencia Política \textbf{·} ITAM \textbf{·} Emiliano Miranda González}
\end{center}
\vspace{-4pt}
\noindent{\color{verdeITAM}\rule{\linewidth}{1pt}}
\vspace{2pt}

\begin{multicols}{3}

% ───────────────────────────── COLUMNA 1 ─────────────────────────────

\tarea{Leer datos}
\begin{itemize}
  \item \co{read\_csv("datos/limpios/x.csv")} --- csv (readr)
  \item \co{read\_dta("x.dta")} --- Stata (haven)
  \item \co{read\_excel("x.xlsx")} --- Excel (readxl)
  \item \co{read\_rds("x.rds")} / \co{write\_rds()}
  \item \co{here::here("datos", "x.csv")} --- ruta al proyecto, nunca \co{setwd()}
  \item \co{janitor::clean\_names()} --- nombres de columna limpios
\end{itemize}

\tarea{Mirar una base}
\begin{itemize}
  \item \co{glimpse(datos)} --- tipos y primeros valores
  \item \co{nrow()} / \co{ncol()} / \co{names()}
  \item \co{distinct(datos, col)} --- valores únicos
  \item \co{sum(is.na(datos\$col))} --- cuántos \co{NA}
\end{itemize}

\tarea{Filtrar y ordenar}
\begin{itemize}
  \item \co{filter(datos, col > 20)}
  \item \co{filter(datos, col \%in\% c(1,2))} --- nunca \co{== 1 | 2}
  \item \co{arrange(datos, desc(col))}
  \item \co{slice\_max(datos, col, n = 5)}
  \item \co{select(datos, a, b, -c)}
\end{itemize}

\tarea{Crear variables}
\begin{itemize}
  \item \co{mutate(datos, ventaja = a - b)}
  \item \co{case\_when(a \~{} "x", TRUE \~{} "y")}
  \item \co{if\_else(cond, si, no)}
  \item Recuerda: \co{mutate()} no guarda solo --- reasigna con \co{<-}
\end{itemize}

\tarea{Agrupar y resumir}
\begin{itemize}
  \item \co{group\_by(datos, region) |>}\newline\co{\ \ summarise(prom = mean(x))}
  \item \co{count(datos, region)} --- atajo de \co{n()}
  \item Media \emph{ponderada} $\neq$ media simple: pondera por tamaño
\end{itemize}

\tarea{Unir}
\begin{itemize}
  \item \co{left\_join(a, b, by = "clave")} --- cruza por clave numérica, no por nombre
  \item \co{anti\_join(a, b, by = "clave")} --- quién quedó sin pareja
  \item Si \co{nrow()} crece tras el cruce: la clave no es única
\end{itemize}

\columnbreak

% ───────────────────────────── COLUMNA 2 ─────────────────────────────

\tarea{Reordenar (pivotear)}
\begin{itemize}
  \item \co{pivot\_longer(datos, cols = a:c,}\newline\co{\ \ names\_to = "v", values\_to = "x")}
  \item \co{pivot\_wider(datos, names\_from = v,}\newline\co{\ \ values\_from = x)}
\end{itemize}

\tarea{Graficar}
\begin{itemize}
  \item \co{ggplot(datos, aes(x, y)) +}
  \item \co{geom\_col()} / \co{geom\_point()} / \co{geom\_line()}
  \item \co{labs(title=, x=, y=, caption=)}
  \item \co{facet\_wrap(\~{}var)} --- un panel por categoría
  \item \co{scale\_fill\_viridis\_c()} / \co{\_d()} --- segura para daltonismo
  \item \co{scale\_fill\_gradient2(midpoint=0)} --- divergente, con cero real
  \item \co{theme\_lab()} --- tema del curso (\co{estilo/tema\_lab.R})
  \item Dentro de \co{aes()}: depende del dato. Fuera: constante.
\end{itemize}

\tarea{Mapear}
\begin{itemize}
  \item \co{st\_read("x.shp")} --- shapefile completo (\co{.shp+.dbf+.shx+.prj})
  \item \co{read\_rds("geo\_simplificado.rds")} --- versión rápida, solo para dibujar
  \item \co{st\_crs(mapa)} --- revisa \emph{siempre} antes de graficar
  \item \co{left\_join(mapa, datos, by="clave")}
  \item \co{geom\_sf(aes(fill=var))}\newline+ \co{theme\_lab\_mapa()}
\end{itemize}

\tarea{Modelar}
\begin{itemize}
  \item \co{lm(y \~{} x, data = datos)} --- el de Estadística I
  \item \co{feols(y \~{} x, data = datos,}\newline\co{\ \ vcov = "hetero")} --- el estándar profesional
  \item \co{modelsummary(list(m1, m2))} --- nunca copies la consola
  \item Coeficiente con unidades: "por cada [x], [y] cambia en promedio..."
\end{itemize}

\tarea{Guardar}
\begin{itemize}
  \item \co{ggsave("figuras/g.pdf", width=7, height=5)}
  \item \co{write\_csv(datos, "x.csv")} / \co{write\_rds()}
  \item \co{dir.create("figuras",}\newline\co{\ \ showWarnings=FALSE)} --- antes de guardar
\end{itemize}

\columnbreak

% ───────────────────────────── COLUMNA 3 ─────────────────────────────

\tarea{Atajos de RStudio que importan}
\begin{itemize}
  \item \co{Ctrl/Cmd + Enter} --- corre la línea o selección
  \item \co{Ctrl/Cmd + Shift + Enter} --- corre todo el script
  \item \co{Alt/Option + "-"} --- inserta \co{<-}
  \item \co{Ctrl/Cmd + Shift + M} --- inserta el pipe (activa el nativo \co{|>} en \emph{Tools > Global Options > Code}, si no, inserta \co{\%>\%})
  \item \co{Tab} --- autocompletar nombres y argumentos
  \item \co{Ctrl/Cmd + S} --- guardar el script
  \item \co{Ctrl/Cmd + Shift + F10} --- reiniciar sesión de R limpia
\end{itemize}

\tarea{Los tres errores más frecuentes}
\begin{itemize}
  \item \co{Error: object 'x' not found} \newline
    No corriste la línea de arriba, o el nombre no coincide exacto (mayúsculas incluidas).
  \item \co{Error in library(tidyverse) :} \newline\co{ there is no package called}\newline
    No lo instalaste. \co{install.packages(} \newline \co{\ \ "tidyverse")} una vez, en la \emph{consola}.
  \item \co{+} parpadeando en la consola \newline
    No es error: falta cerrar un paréntesis o una comilla. \co{Esc} y revisa.
\end{itemize}

\vspace{4pt}
\noindent{\color{grisTenue}\rule{\linewidth}{0.4pt}}\\[2pt]
{\fontsize{7.6}{8.6}\selectfont\color{grisTexto}
Pipe nativo \co{|>} siempre, nunca \co{setwd()}. Catálogo completo de errores en
\co{recursos/errores-comunes.qmd}; inglés técnico traducido en \co{recursos/glosario.qmd}.
}

\end{multicols}

\end{document}
